\subsection{Fatoriais e binomiais}

\prob{1}{fatbinom1} Como calcular $C(n, k) = {n \choose k} = \dfrac{n!}{k!(n-k)!}$?

\textbf{Solução:} Você pode calcular cada fatorial. Você pode calcular o produto a partir de um certo ponto no numerador e depois dividir pelo fatorial que sobrou. Você pode utilizar a recorrência $C(n, k) = C(n-1, k-1) + C(n-1, k)$. Você pode ter os fatoriais precomputados e calcular em $O(1)$.

\textbf{Propriedades:}
$$ {n \choose k} = {n \choose n -k}$$
$${n \choose k} = \dfrac{n}{k} {n-1 \choose k-1}$$
$${n \choose k} = {n-1 \choose k-1} + {n-1 \choose k}$$
$$\sum\limits_{k=0}^n {n \choose k} = 2^n$$
$$\sum\limits_{m=0}^n {m \choose k} = {n + 1 \choose k+1}$$
$$\sum\limits_{k=0}^m {n+k \choose k} = {n+m+1 \choose m}$$
$$\sum\limits_{k=0}^n {n \choose k}^2 = {2n \choose n}$$
$$\sum\limits_{k=0}^n k {n \choose k} = n \cdot 2^{n-1}$$
$$\sum\limits_{k=0}^n {n-k \choose k} = F_{n+1}$$
$$(a+b)^n = \sum\limits_{k=0}^n {n \choose k} a^k b^{n-k}$$
$${m + n \choose r} = \sum\limits_{k=0}^r {m \choose k} {n \choose r - k}$$
$$\sum\limits_{k=0}^m (-1)^m {n \choose k} = (-1)^m {n-1 \choose m}$$

\textbf{Obs.:} Se o "denominador" do coeficiente binominal for maior que o "numerador", definimo-lo como zero.

\prob{2}{fatbinom2} Calcular $C(n, k) \mod m$.

\textbf{Solução:} fatorar $m$ em potências de primo e juntar tudo com o teorema chinês dos restos. Para achar $C(n, k) \mod p^e$, seja $c(x)$ o maior inteiro tal que $p^{c(x)} | x!$. Seja $g(x) = \dfrac{x!}{p^{c(x)}}$. Você pode calcular os $g$ e os $c$ com dp por exemplo. Aí no binomial você teria que lidar com

$$ \dfrac{g(n)}{g(k)g(n-k)} p^{c(n) - c(k) - c(n-k)}$$

Como os $g$ agora são coprimos com $p$ dá pra calcular de boa o mod deles por potência de primo. E, quanto à outra parte, é só multiplicar no final.

\prob{3}{fatbinom3} $C(n, k) \mod m$ pra $n$ grande e $m$ pequeno.

\textbf{Solução 1 (m primo):} Usa as técnicas do problema $2$, que leva $O(m \log_m n)$.

\textbf{Solução 2 (m primo):} Usa o teorema de Lucas. Esse teorema diz que se $n = n_l p^l + \dots + n_0$ e $k = k_l p^l + \dots + k_0$ forem a expressão de $n$ e $k$ na base $p$, então
$$ {n \choose k} \equiv \prod\limits_{i=0}^l {n_i \choose k_i} \mod p$$

\textbf{Solução 3 (m livre de quadrados, fatores pequenos):} Usa o teorema de lucas e combina tudo com o teorema chinês dos restos.

\textbf{Solução 4 (m qualquer):} Usa o teorema de lucas generalizado e junta com o teorema chinês dos restos. Seja $(k!)_p$ o produto de todos os inteiros menores ou iguais a $k$ não divisíveis por $p$. Suponha que queremos achar ${n \choose m}$ (com $r = n - m \geq 0$) módulo $p^f$. Escreva $n = n_0 + n_1 p + \dots + n_s p^s$ em base $p$. Seja $N_j$ o menor resíduo positivo de $\lfloor n/p^j \rfloor ( \mod p^f)$ (de modo que $N_j = n_j + \dots + n_{j+f-1} p^{f-1}$). Defina de maneira análoga $m_j, M_j, r_j, R_j$. Seja $e_j$ o número de índices $i \geq j$ tais que $n_i < m_i$. Então:

$$\dfrac{1}{p^{e_0}} {n \choose m} \equiv (\pm 1)^{e_{f-1}} \prod\limits_{i=0}^s \dfrac{(N_i!)_p}{(M_i!)_p (R_i!)_p} ( \mod p^f)$$

onde o $\pm 1$ é $-1$, exceto no caso $p=2$ e $f \geq 3$.

\prob{4}{fatbinom4} Você quer ir de $(a,b)$ para $(c, d)$ indo apenas pra cima/baixo/esq/dir. Qual o número de caminhos com a menor distância?

\textbf{Solução:} Seja $\Delta_x = |a-c|$ e $\Delta_y = |b-d|$. Então, a resposta é ${\Delta_x + \Delta_y \choose \Delta_x}$.

\prob{5}{fatbinom5} Você tem uma array de $n=10^5$ e vai processar $m=10^5$ queries que adicionam ${k \choose k}$, ${k+1 \choose k}, \dots$ no intervalo $[l, r]$. Sabendo que $k \leq 100$, dizer como a array fica no final.

\textbf{Solução:} Faça mais $101$ arrays de "subtrações" de prefixo. Diga que $vet[0][i] = a[i] - a[i-1]$ e $vet[i][j] = vet[i-1][j] - vet[i-1][j-1]$. Se uma query envolve um $k$ arbitrário, incremente $vet[k][l]$ e decremente $vet[k][r+1]$. Depois de em todas as queries fazer isso, para cada $k$, reconstrua a array original a partir de $vet[k]$. Você vai obter várias arrays modificadas. Cada uma delas considera todas as queries para determinado $k$ agindo no vetor. Então você soma todas elas e ajusta.

\prob{6}{fatbinom6} Você quer saber para quantos/quais $k$ o número ${n \choose k}$ é ímpar.

\textbf{Solução:} Pelo teorema de Lucas, para ser ímpar, a cada bit 1 de $k$, não pode corresponder um bit $0$ de $n$. Em outras palavras, cada bit $0$ de $n$ corresponderá a um bit $0$ de k. E pra cada bit $1$ de $n$, o bit correspondente de $k$ não importa. Daí, você pode iterar de $0$ até $n$ e pegar os $k$ que satisfazem o que queremos com isso. E para pegar a quantidade é só $2^{n_1}$ onde $n_1$ é o número de bits $1$ da representação decimal de $n$.

\prob{7}{fatbinom7} Dada uma string de parêntesis, como determinar o número de substrings (não necessariamente contíguas) iguais a (), (()), ((())), etc?

\textbf{Solução:} Para a string (((...)))) com $x$ ( e $y$ ), a resposta é ${x+y \choose x}$ (por ex via vandermonde). Agora, itere pelos ( da string e para cada um deles, conte o número de substrings em que este é o último opening bracket. Para isso, suponha que tem $x$ opening brackets (estritamente) antes dele e $y$ closing brackets depois. Aí como esse ( atual deve estar incluído na resposta, então queremos contar as substrings de (((...))). Então ponha uma lista no último opening bracket e vá somando as escolhas de $r-1$ opening brackets e $r$ closing brackets e some. Pela mesma lógica de vandermonde, agora a respostá é ${x+y \choose x+1}$.

\prob{8}{fatbinom8} Tem uma string s de $n \leq 10^6$ letras de a ate z. Cada conjunto de letras contíguas é uma "colônia" da especie de bacteria correspondente a essa letra (então número de espécies é menor ou igual a 26 e pode ter mais de uma colônia de uma dada espécie). A cada instante ocorre um ataque. Um ataque é tipo "aabb" vira "aaab" e é sempre um caractere. A pergunta é o número de possíveis estados que podem ocorrer. Por exemplo, se for "aab", temos "aab", "aaa", "abb", "bbb", então a resposta é $5$;

\textbf{Solução:} Converta cada colonia para um caractere por ex "aaaabbccc" vira "abc". É óbvio que uma string t de $n$ caracteres é um estado válido se, e só se, ao eu converter cada colonia dela pra um caractere, isso é uma substring do "abc". Ou seja, dada uma substring de "abc", basta eu ver quantas strings t são tais que a compressão dessa string t seja essa substring. Digamos que essa substring tenha $m$ caracteres. É fácil contar as strings t porque são o número de soluções de $x_1 + \dots + x_m = n$ com $x_i \geq 1$. Ou seja, ${n - 1 \choose m-1}$. Então temos que contar o número de substrings (não necessariamente contíguas) de "abc" (isto é, da compressão de s) que não possuem dois caractéres consecutivos iguais e que tenham "m" caracteres para cada valor de "m". Dá pra fazer isso com uma dp[c][l] - numero de substrings de tamanho l que terminam em c. dp[ci][l] = 1 + soma[c != ci] dp[c, l-1]

\prob{9}{fatbinom9} Tem $n$ pontos com coordenadas inteiras no plano. Pra cada ponto, ou você põe uma reta horizontal passandro por ele, ou uma vertical, ou você não faz nada. Dadas as coordenadas, qual o número de desenhos?

\textbf{Solução:} Imagina que você tem um grafo em que os pontos são vértices e tem uma aresta se e só se tiverem o mesmo x ou o mesmo y. A resposta é claramente o produto das respostas pra cada componente. Em uma dada componente, se nela não tiver nenhum "ciclo" (não ciclo, mas 4 pontos que formam um retangulo paralelo aos eixos coordenados), então a resposta pra ela é $2^{x+y}-1$. Se tiver, é $2^{x+y}$, em que $x$ é a quantidade de coordenadas $x$ distintas e $y$ analogamente.

\prob{10}{fatbinom10} Palitos e bolinhas. Você tem $x_i \geq 0$ e quer saber o número de soluções de $x_1 + \dots + x_n = N$.

\textbf{Solução:} ${N + n - 1 \choose n-1}$

\prob{11}{fatbinom11} Palitos e bolinhas estendido. Você tem $0 \leq x_i \leq r_i$ e quer saber o número de soluções de $x_1 + \dots + x_n = N$.

\textbf{Solução:} Argumenta por inclusão e exclusão no palitos e bolinhas normal. Vai chegar em
$$ \sum\limits_{S \subseteq \{1,2, \dots,n \}} (-1)^{|S|} {N + n - 1 - \sum\limits_{i \in S} r_i + 1 \choose n-1}$$
No caso particular em que $r_i = r$, temos
$$ \sum\limits_{k=0}^{\left\lfloor N/(r+1) \right\rfloor} (-1)^k {n \choose k} {N - k (r+1) + n - 1 \choose n -1}$$

\prob{12}{fatbinom12} Você tem um sanduíche de $N \leq 10^{18}$ metros e quer cortar no menor número de pedaços com comprimento positivo que é tal que é possível que cada pedaço fique com tamanho $\leq K \leq 10^{18}$ mod $M \leq 10^6$

\textbf{Solução:} O número de pedaços $q$ satisfaz $\lceil N/q \rceil \leq K$. Na verdade queremos o menor $q$ que satisfaz isso. Esse $q$ é $\lceil N / K \rceil$. Se o tamanho do sanduíche fosse $qK$, a resposta seria $1$. Mas o tamanho dele é $N$. Note que $qK - N $, ou seja, estamos descendo do menor múltiplo de $K$ maior ou igual a $N$ para o próprio $N$. Então essa diminuição é menor que $K$. Ou seja, basta resolvermos $x_1 + \dots + x_q = qK - N$ em que $x_i < K$ (mas esta é uma restrição desnecessária, visto que $qK - N < K$), onde $x_i$ representa o quanto tiramos do sanduíche $i$. Por isso, a resposta final é ${qK - N + q - 1 \choose q - 1}$. Para calcular isso, temos binomiais muito grandes e arbitrários. Agora fatora o $M$ e usa o teorema de lucas generalizado em cada potência de primo e junta tudo com o teorema chinês dos restos.